\documentclass[letterpaper,10pt]{article}
\usepackage{amsmath}
\usepackage{amssymb}
\usepackage{amsthm}
\usepackage[margin=1in]{geometry}
\usepackage{enumitem}
\usepackage{graphicx} % Required for including images

% Custom theorem-like environments
\theoremstyle{definition}
\newtheorem{definition}{Definition}

\theoremstyle{plain}
\newtheorem{theorem}{Theorem}

\begin{document}

\title{CSE 291: Unsupervised learning \\ Spring 2008}
\author{}
\date{Lecture 5 Finding meaningful clusters in data}
\maketitle

So far we've been in the vector quantization mindset, where we want to approximate a data set by a small number of representatives, and the quality of the approximation is measured by a precise distortion function. Now we move to a different use of clustering: to discover interesting structure in data.

This immediately seems fuzzy and ill-defined, and it is unclear how to make a principled choice among the plethora of clustering heuristics that present themselves. But if we step back and take a slightly more abstract view, some interesting conclusions and directions emerge.

We'll assume that we have a data set S, along with some notion of interpoint distances. For the latter, we use the following formalism.

\begin{definition}
$d:S\times S\rightarrow\mathbb{R}$ is a distance function if it satisfies the following two properties:
\begin{enumerate}
    \item $d(x,y)\ge0$, with equality if and only if $x=y$
    \item $d(x,y)=d(y,x).$
\end{enumerate}
In particular, any metric is a distance function, although a distance function need not satisfy the triangle inequality.
\end{definition}

Given only S and d, what is a reasonable way to generate a clustering?

\section{Kleinberg's axiomatic framework for clustering}
Kleinberg formalizes a clustering function as a mapping that takes as input a distance function d on n points and returns a partition $\Gamma$ of $[n]=\{1,2,...,n\}$; if this mapping is f, we will write $f(d)=\Gamma$. For instance,
\[
f\left(\begin{pmatrix}0&10&10\\ 10&0&1\\ 10&1&0\end{pmatrix}\right)=(\{1\},\{2,3\})
\]
sends a distance function on three points to a clustering with two clusters. Notice that the clustering function has to choose the number of clusters.

\subsection{Axioms for reasonable clustering functions}
What properties might one expect a good clustering function to possess? Here are some possibilities.
\begin{itemize}
    \item \textbf{Scale invariance.} For any $\alpha>0$ and any d, $f(\alpha\cdot d)=f(d)$.
    \item \textbf{Richness.} $Range(f)=$ \{all possible partitions of $[n]\}$.
    \item \textbf{Consistency.} If $f(d)=\Gamma$ and $d^{\prime}$ is a $\Gamma$-enhancing transformation of d, then $f(d^{\prime})=\Gamma.$
\end{itemize}
For the last property, a $\Gamma$-enhancing transformation is one that makes a distance function even more suggestive of the clustering $\Gamma$. More precisely:

\begin{definition}
$d^{\prime}$ is a $\Gamma$-enhancing transformation of d if
\begin{itemize}
    \item $d^{\prime}(i,j)\le d(i,j)$ for i, j in the same cluster of $\Gamma$
    \item $d^{\prime}(i,j)\ge d(i,j)$ for i, j in different clusters of $\Gamma$
\end{itemize}
\end{definition}

For instance, in the figure below, the left side shows a set of six points (with Euclidean distances d) and a particular clustering $\Gamma$ consisting of three clusters. The right side shows a modified configuration whose distance function $d^{\prime}$ is a $\Gamma$-transformation of d.

%\begin{figure}[h!]
%    \centering
%    \includegraphics[width=0.8\textwidth]{lec5_fig1.png} % Placeholder for the image
%    \caption{Example of $\Gamma$-enhancing transformation. Left: Original configuration with distance d and clustering $\Gamma$. Right: Modified configuration with distance $d'$.}
%\end{figure}

What kinds of algorithms satisfy the axioms above? Here are three illustrative examples.
\begin{enumerate}
    \item Run single linkage till you get k clusters. This satisfies scale invariance and consistency, but not richness.
    \item Run single linkage till distances exceed $c\cdot \max_{i,j}d(i,j)$, where c is some constant. This satisfies scale invariance and richness but not consistency.
    \item Run single linkage until distances exceed some threshold r. This satisfies richness and consistency but not scale invariance.
\end{enumerate}
In fact, Kleinberg shows that there is no clustering function that satisfies all three axioms.

\subsection{An impossibility result}
The basic problem is that there is no consistent way to choose a level of granularity. Does the following picture contain 2 clusters, or 3, or 12? Let's see how this problem makes it impossible to satisfy all three axioms. Suppose we have eight points whose interpoint distances d look like this:

%\begin{figure}[h!]
%    \centering
%    \includegraphics[width=0.8\textwidth]{lec5_fig2.png} % Placeholder for the image
%    \caption{Example of eight points with interpoint distances d.}
%\end{figure}

This looks like one cluster, so suppose $f(d)=\Gamma$ consists of a single cluster. Then a $\Gamma$-enhancing transformation yields the following $d^{\prime}$ in which the four points initially on the left are now almost on top of each other, and likewise with the four points initially on the right.

%\begin{figure}[h!]
%    \centering
%    \includegraphics[width=0.8\textwidth]{lec5_fig3.png} % Placeholder for the image
%    \caption{Example of $\Gamma$-enhancing transformation yielding $d'$.}
%\end{figure}

Rescaling, we get $d^{\prime\prime}:$

%\begin{figure}[h!]
%    \centering
%    \includegraphics[width=0.8\textwidth]{lec5_fig4.png} % Placeholder for the image
%    \caption{Rescaled version $d''$ of $d'$.}
%\end{figure}

Now, consistency and scale invariance would enforce $f(d^{\prime\prime})=f(d^{\prime})=f(d)$. But d looks like one cluster while $d^{\prime\prime}$ looks like two! To make this intuition more precise, we introduce a key concept.

\begin{definition}
Partition $\Gamma^{\prime}$ is a refinement of $\Gamma$ (written $\Gamma^{\prime}\le\Gamma$) if for every $C^{\prime}\in\Gamma^{\prime}$ there is some $C\in\Gamma$ such that $C^{\prime}\subset C$.
\end{definition}
In other words, $\Gamma$ can be obtained by merging clusters of $\Gamma^{\prime}$. For instance:

%\begin{figure}[h!]
%    \centering
%    \includegraphics[width=0.6\textwidth]{lec5_fig5.png} % Placeholder for the image
%    \caption{Example of partition refinement.}
%\end{figure}

\begin{theorem}
Pick any $\Gamma_{0}\le\Gamma_{1}$ with $\Gamma_{0}\ne\Gamma_{1}$. If f satisfies scale invariance and consistency then $\Gamma_{0}$ and $\Gamma_{1}$ cannot both be in $range(f)$.
\end{theorem}
\begin{proof}
Suppose, on the contrary, that $f(d_{0})=\Gamma_{0}$ and $f(d_{1})=\Gamma_{1}$. We'll arrive at a contradiction.
%\begin{figure}[h!]
%    \centering
%    \includegraphics[width=0.6\textwidth]{lec5_fig6.png} % Placeholder for the image
%    \caption{Diagram illustrating assumptions for proof of Theorem 4.}
%\end{figure}
Let $a_{0}=\min\{d_{0}(i,j):i\ne j\}$ and $b_{0}=\max\{d_{0}(i,j)\}$; likewise, let $a_{1}$ and $b_{1}$ be the smallest and largest values of $d_{1}$. Now construct a new distance function d:
\[
d(i,j) =
\begin{cases}
    \alpha_{0} & \text{if } i, j \text{ in same cluster of } \Gamma_{0} \\
    a_{1} & \text{if } i, j \text{ in same cluster of } \Gamma_{1} \text{ but not } \Gamma_{0} \\
    b_{1} & \text{if } i, j \text{ in different clusters of } \Gamma_{1}
\end{cases}
\]
where $\alpha_{0}$ is chosen such that $0 < \alpha_{0} < a_{1}$.

Distance function d is a $\Gamma_{1}$-enhancing transformation of $d_{1}$. Therefore, $f(d)=f(d_{1})=\Gamma_{1}$. But now consider distance function $d^{\prime}=(b_{0}/a_{1})\cdot d$; notice that
\[
d^{\prime}(i,j)=\begin{cases}
    a_{0} & \text{if } i,j \text{ in same cluster of } \Gamma_{0} \\
    \ge b_{0} & \text{if } i,j \text{ in different clusters of } \Gamma_{0}
\end{cases}
\]
Therefore $d^{\prime}$ is a $\Gamma_{0}$-enhancing transformation of $d_{0}$ and $f(d^{\prime})=f(d_{0})=\Gamma_{0}$. But d' is also a scaled version of d, so we need $f(d)=f(d^{\prime})$, a contradiction.
\end{proof}

\subsubsection*{Problem 1}
Develop an axiomatic framework for hierarchical clustering. This will circumvent the difficulty of automatically choosing a level of granularity.

\section{Finding prominent clusterings}
In some situations, we primarily want a clustering algorithm to find prominent groups in the data, if such groups exist. If, on the other hand, the data has no well-defined clusters, then it doesn't matter what the algorithm does.

\end{document}
